\documentclass[11pt]{article}
\usepackage{amsmath,amssymb,graphicx,hyperref,booktabs}
\usepackage[margin=1in]{geometry}
\title{Q-Trust: Verifiable Post-Quantum Migration \\ via CBOM Discovery, GNN Planning, and On-Chain Attestation}
\author{Q-Trust Contributors \\ \texttt{humoge7502/q-trust}}
\date{September 2026 (draft)}
\begin{document}
\maketitle
\begin{abstract}
NIST IR~8547 schedules RSA/ECC disallowance for 2030--2035, yet most
organizations cannot enumerate their own cryptographic estate, let alone plan
its migration. Q-Trust is an end-to-end open system: (1)~a multi-language
scanner emitting CycloneDX-native CBOMs with NIST/CNSA scoring; (2)~a
CodeBERTa discovery model (F1~0.9525, repo-disjoint held-out); (3)~a GNN/RL
migration planner evaluated by 42-fold leave-one-out on 42 host-disjoint
enterprise CBOMs (277 hosts, zero overlap); (4)~EIP-712 gasless attestation
anchored on Base L2 behind a timelock. Our headline finding is a tie: the GNN
($\tau_b$~0.7168) reproduces the doctrine heuristic ($\tau_b$~0.7210) with
0~wins / 41~ties / 1~loss. We publish the tie deliberately: leakage audits,
deterministic splits, and honest-failure harnesses show the doctrine is
near-optimal on today's estates, and identify the missing ingredient
(blinded expert pairwise labels). All artifacts carry config, seed, and
metrics; all demo data is manifest-labeled \texttt{SYNTHETIC\_DEMO}.
\end{abstract}

\section{Introduction}
``Harvest now, decrypt later'' makes post-quantum migration a scheduling
problem under a hard deadline. Migration requires three capabilities rarely
found together: \emph{inventory} (what crypto exists where), \emph{ordering}
(what migrates first under constrained budgets), and \emph{proof} (that the
migration happened and has not regressed). Q-Trust implements all three in
one open pipeline: scanner $\to$ planner $\to$ on-chain attestation.

\section{Related work}
Static crypto-misuse detectors (CryptoGuard, CogniCrypt) solve a different
task than ours: they detect \emph{misuse} given known crypto; Q-Trust solves
\emph{usage discovery} across polyglot estates. Measured on CryptoAPI-Bench
file usage, our deterministic layers reach usage recall 0.877 at file-level
misuse precision 0.932; on the Apache corpus (71/121 mapped jars) F1~0.688.
We claim no superiority over misuse detectors: the tasks differ, and our
report states this framing explicitly.

\section{Discovery}
A CodeBERTa-small-v1 fine-tune (4 epochs, CUDA, seed~42) over 12,462 real
files (6,636 crypto-labeled, 34 repositories) reaches P~0.952 / R~0.953 /
F1~0.9525 on a repo-disjoint held-out set ($n{=}2{,}415$). The split
reproduces bit-identically; cross-repo exact-duplicate contamination is 0.4\%,
bounded far below the +0.11 F1 gain over the best baseline.

\section{Planning and the honest tie}
The GNN ranks migration order per estate. Evaluated by 42-fold LOO over
42~host-disjoint CBOMs: $\tau_b$~0.7168 vs doctrine heuristic 0.7210
($\Delta$~$-0.0042$, identical medians), top-10 accuracy 1.0 on all folds.
An RL agent independently ties (140.34~$\pm$~8.13 vs 140.62). Two independent
methods converging to the heuristic is evidence the doctrine is near-optimal
given current observables --- not evidence of model failure. The corpus ships
a schema-complete RiskBench (10,000 pairs) whose labels are synthetic demo
data until 5--10 blinded experts annotate v1.

\section{Side channels and attestation}
A trace-anomaly detector retrained on 54 real liboqs timing sets verifies
46~clean + 8~low-risk traces and flags 51/54 leak-injected sets (misses:
FALCON-512/1024 and SLH-DSA-256s keygen at $n{\le}200$, recorded openly).
Migration evidence is anchored via EIP-712 gasless attestations to UUPS
registries behind a 7-day timelock; 213 Forge tests (incl.\ adversarial and
invariant suites) and a clean Slither run cover project sources.

\section{Evaluation integrity}
Every headline metric traces to an artifact (\texttt{docs/TRUTH\_AUDIT.md}
classifies each as REAL/SYNTHETIC/DEMO). Unwired harnesses report
\texttt{not\_available} rather than fabricating values; the eval entry point
reports \texttt{is\_demo: false} only over the verified real splits.

\section{Limitations}
No external contract audit yet (dossier prepared); PQC ACVP vectors are
fixtures, not conformance claims; trace coverage is 5 of 10+ target parameter
sets; enterprise scale ($\geq$1K org CBOMs) awaits volunteer scans.

\section{Conclusion}
Q-Trust demonstrates that honest evaluation is a feature: a leakage-free tie
is more useful than a leaky win, because it prescribes the next experiment.
Code, data manifests, and the 60-second reproduction path are public.

\begin{thebibliography}{9}
\bibitem{nist8547} NIST IR 8547, Transition to Post-Quantum Cryptography.
\bibitem{fips203} NIST FIPS 203, ML-KEM. \bibitem{fips204} NIST FIPS 204, ML-DSA.
\bibitem{cyclonedx} OWASP CycloneDX 1.7, CBOM specification.
\bibitem{cryptobench} CryptoAPI-Bench; ApacheCryptoAPI-Bench corpora.
\bibitem{codeberta} HuggingFace CodeBERTa-small-v1.
\end{thebibliography}
\end{document}
